\section{Glossaire de lecture}\label{sec:glossaire}
Les entrées suivent l'ordre de première apparition dans le corps du
texte, après le résumé. Les symboles introduits ensemble sont regroupés ;
leurs définitions opératoires restent celles des sections correspondantes.

\begin{small}
\begin{longtable}{@{}>{\raggedright\arraybackslash}p{.27\linewidth}p{.66\linewidth}@{}}
\toprule Terme ou notation & Sens dans cet article \\
\midrule\endhead
Effet rebond & Écart à l'économie de consommation attendue après un gain
d'efficacité. Ici, l'expérience mesure une réponse productive non
proportionnelle, sans calibration d'une consommation finale. \\
Endogène ; émergence & Produit par les interactions et l'évolution du
système, et non imposé comme cible macroscopique. \\
Auto-organisation critique (SOC) & Hypothèse d'un régime se maintenant
spontanément près d'un seuil critique. Elle n'est pas démontrée ici. \\
Entité & Agent générique producteur, susceptible de prêter, d'emprunter
et de disparaître ; aucun de ces rôles n'est prédéfini. \\
$K_i$ ; stock productif & Capital physique mobilisable de l'entité $i$,
en joules du modèle ; distinct de sa valeur nette. \\
Technologie $f_i(K)=A_iK^{\gamma_i}$ & Production ajoutée au stock
pendant un pas ; $A_i$ est le coefficient d'extraction, $\gamma_i$
l'exposant de production. \\
$K^\circ$, $\gamma^\circ$ & Dotation de naissance et exposant de
production à la naissance dans le régime de référence homogène. \\
Principal $q$ ; service $rq$ & Montant nominal d'un prêt et intérêts
exigibles par pas ; le principal n'est pas amorti. \\
$C_i$, $D_i$, $NW_i$ & Créances, dettes et valeur nette
(\emph{net worth}) : $NW_i=K_i+C_i-D_i$. \\
$\lambda$, $\delta$, $\sigma$, $\rho$ & Intensité de naissance,
fraction dépréciée par pas, amplitude du choc log-normal et intensité
des rencontres de marché. \\
Institution de crédit & Règles fixant le sens, le montant et le taux
des prêts ; elle ne désigne pas ici une banque particulière. \\
Invariant ; bilan réel & Relation imposée par les opérations du modèle.
Le bilan réel suit les apports, la production et les pertes du stock. \\
$B_t$, $\eta_i$ & Nombre de naissances au pas $t$ et choc logarithmique
individuel, renouvelé à chaque pas. \\
Concavité & Diminution du rendement marginal avec le stock, pour
$0<\gamma<1$ ; origine du gain possible de réallocation. \\
Défaut de service & Incapacité à payer intégralement les intérêts dus
au pas considéré. \\
Stock ; flux ; $\Delta t$ & Quantité détenue ; quantité transférée ou
produite pendant un pas ; durée physique éventuelle de ce pas. \\
Exergie & Capacité d'une ressource à fournir du travail relativement à
un environnement de référence ; sa correspondance avec les droits
économiques reste ici une hypothèse. \\
$K_{\rm aut}$, $x_i$ & Stock fixe déterministe d'une entité isolée
sans bruit ni crédit ; capital réduit $x_i=K_i/K_{\rm aut}$. \\
$\kappa^\circ$, $a$ & Dotation réduite $K^\circ/K_{\rm aut}$ ;
coefficient alternatif $a=A(K^\circ)^{\gamma-1}$. \\
Rééchelonnement ; compensation & Le premier transforme tout l'état et
ses unités ; la seconde adapte ici la dotation des futures entrantes
pour conserver $\kappa^\circ$ après la hausse de $A$. \\
Tension $T$ ; $K_{\rm eq}$ & Rapport $K_{\rm aut}/K_{\rm eq}$ ;
stock fictif donnant la production moyenne par productrice.
Dans les protocoles, $T$ peut aussi désigner l'horizon d'une simulation. \\
$Y$, $N_{\rm prod}$ & Production totale du pas et nombre d'entités
présentes à la phase productive. \\
Rendement marginal $f'(K)$ & Production supplémentaire à la marge
par unité de stock supplémentaire. \\
$L$, $\Delta$, $p$ & Perte productive de la prêteuse, surplus productif
joint et part de ce surplus attribuée à la prêteuse : $rq=L+p\Delta$. \\
Faillite ; cascade ; avalanche & Suppression d'une entité en défaut
ou insolvable ; propagation des pertes ; composante faiblement connexe
des pertes entre mortes du même pas. \\
Racine ; génération ; $b_1$, $b_2$ & Mort initiale ; rang de propagation
intra-pas ; fraction des mortes non racines et nombre moyen de liens
causals entre générations successives par morte. \\
Graine ; appariement ; IC95 & Initialisation du générateur aléatoire ;
branches issues d'un même état ; intervalle de confiance à 95\,\%
calculé entre répétitions, sauf indication contraire. \\
$k$ ; stationnarité & Taille historique du bassin de rencontre, fixée
à deux dans le protocole principal ; invariance temporelle des propriétés
statistiques, non démontrée par un seul test de dérive. \\
Loi de Little & Relation entre effectif moyen, flux d'entrée et durée
de vie moyenne sous conditions de stabilité et de maîtrise des bords. \\
$S$, $\alpha$, $S_c$, $Z$ & Taille d'avalanche, exposant, coupure
exponentielle et constante de normalisation de sa loi ajustée. \\
Bootstrap ; $R^2$ & Rééchantillonnage de réalisations entières ;
score descriptif d'ajustement, non preuve d'une famille de lois. \\
Contraction ; $D$, $\beta$ & Suite de pas de baisse d'une grandeur ;
durée de cette suite et taux de décroissance de sa survie exponentielle.
$D$ sans indice dans ce contexte n'est pas une dette. \\
Gini ; Lorenz & Indice synthétique d'inégalité ; courbe de part cumulée
d'une grandeur en fonction de la part d'entités classées par elle. \\
Queue ; Pareto ; log-normale & Domaine des grandes valeurs ; famille
à décroissance en puissance ; famille dont le logarithme est normal. \\
Élasticité finie $\varepsilon$ & Rapport du logarithme de la réponse
productive au logarithme du facteur technique $m$, ici $m=1{,}5$. \\
Quantile ; décile ; cohorte & Seuil de classement ; groupe de 10\,\%
reclassé à chaque date ; ensemble d'individus suivis par identité. \\
Covariance & Mesure de variation conjointe utilisée dans l'identité
de pondération du partage ; elle ne démontre pas une causalité. \\
Taux de charge & Intérêts versés ou dus, selon la définition,
rapportés à la production plus les intérêts reçus ; distinct d'un
rapport stock de dette/production. \\
Vuong & Comparaison de vraisemblances de modèles sur un même
échantillon ; ses hypothèses et ses limites sont explicitées en annexe. \\
\bottomrule
\end{longtable}
\end{small}
